% ==============================================================================
% T0 Theory / FFGFT: Document 281 (Chapter Content, EN)
% Title: The Golden Ratio in FFGFT -- Notation, Roles and Status
% Author: Johann Pascher
% Date: 15 June 2026
% Purpose: A consolidating clarification. The earlier documents are NOT
%          changed; this document fixes the canonical reading and points to
%          all substantial occurrences of phi in the corpus.
% References: Doc. 005, 009, 018, 019, 022, 023, 028, 035, 036, 083, 133,
%          147, 149, 150, 153, 154, 159, 172, 187, 188, 190 (P35/P37), 271,
%          272, 274, 275; mass derivation Doc. 006.
% ==============================================================================

\section*{Abstract}
The golden ratio $\varphi = (1+\sqrt5)/2$ appears in the FFGFT corpus in
several clearly distinguishable roles -- and under two different glyphs that
partly collide with other quantities (scalar field, angle, potential, wave
function). This document does \emph{not} change the earlier texts; it fixes the
canonical notation and reading, assigns each substantial occurrence to a role,
and cleanly separates what is \emph{derived} from $\varphi$, what is merely
\emph{asserted}, and what is simply \emph{miscalculated}. Key statements:
(i)~$\varphi$ is nowhere a foundation in FFGFT -- that is $\xi=4/30000$;
(ii)~the only $\varphi$ usage with genuine geometric content is
$\arctan(\varphi^{-3})$ in the CP phase (Doc.~172), because
$\varphi^{-3}=\sqrt5-2$ is the natural intermediate scale of Fibonacci
convergence; (iii)~$\varphi$ mass ladders (Doc.~172/188) are heuristic,
numerically erroneous and fall under the $\varphi^N$ triviality warning P35;
(iv)~$1/\varphi$ is a transit point, not a fixed point (Doc.~275);
(v)~$\pi$ and $e$ are nowhere derived.

% ==============================================================================
\section{Occasion and Scope}
% ==============================================================================

While maintaining the corpus a notation problem became apparent: the same
symbol \verb|\phi| means different things in different documents, and the
golden ratio itself is written sometimes as $\phi$ (closed form, from
\verb|\phi|), sometimes as $\varphi$ (open form, from \verb|\varphi|). Since
the old documents are finalised in version v1.1.2 and are no longer rewritten
retroactively, this document serves as a \emph{binding reference}: anyone who
encounters $\varphi$ or $\phi$ in an older text looks up here which role is
meant.

\begin{important}
\textbf{Scope rule.} This document is normative for the \emph{reading}, not for
the \emph{wording} of the old documents. Where an old text and this document
contradict each other, the version recorded here -- reconciled against
Doc.~190 (P35/P37) and Doc.~275 -- governs the physical content.
\end{important}

% ==============================================================================
\section{Canonical Notation}
% ==============================================================================

The following assignment is the canonical convention. It also describes
\emph{how} the corpus actually uses the symbols -- with the known
inconsistency for the golden ratio.

\begin{center}
\begin{adjustbox}{max width=\textwidth}
\begin{tabular}{l l l l}
\hline
\textbf{Command} & \textbf{Glyph} & \textbf{Canonical meaning} & \textbf{Remark} \\
\hline
\verb|\varphi| & $\varphi$ (open) & golden ratio $(1+\sqrt5)/2$ & recommended standard symbol \\
\verb|\phi|    & $\phi$ (closed) & scalar field / angle / index & NOT the golden ratio \\
\verb|\Phi|    & $\Phi$ (capital) & potential / complex field & e.g.\ $\Phi=\rho e^{i\theta}$ \\
\verb|\psi|    & $\psi$ & wave function & quantum mechanics \\
\hline
\end{tabular}
\end{adjustbox}
\end{center}

\begin{caution}
\textbf{Known legacy issue.} In v1.1.2 several documents write the
\emph{golden ratio} erroneously as \verb|\phi| instead of \verb|\varphi|
(among others Doc.~005, 009, 022, 023, 028, 035, 036, 083, 133, 147, 187,
188). At the same time, field-theory documents use \verb|\phi| correctly for
the \emph{scalar field} (e.g.\ Doc.~202: $\phi_{n,l,j}$ with
$\Box\,\phi_{n,l,j}=-m^2\phi_{n,l,j}$) and \verb|\phi| as a
\emph{winding index} $n_\phi$ in $(n_\theta,n_\phi)$ (e.g.\ Doc.~051, 202).
Within a \emph{single} document $\phi$ is unambiguous in each case; the
overloading exists only \emph{between} documents. Hence one must not replace
\verb|\phi|$\to$\verb|\varphi| wholesale -- that would destroy the
field/angle documents.
\end{caution}

\noindent\textbf{Numerical convention.} Throughout,
\begin{equation}
  \varphi = \frac{1+\sqrt5}{2} = 1.6180339887\ldots, \qquad
  \varphi^2 = \varphi + 1, \qquad
  \varphi^{-1} = \varphi - 1 = 0.6180\ldots,
\end{equation}
\begin{equation}
  \varphi^{-3} = \sqrt5 - 2 = 0.2360679\ldots, \qquad
  \arctan(\varphi^{-3}) = 13.2825^\circ .
\end{equation}

% ==============================================================================
\section{The Roles of $\varphi$ in the Corpus}
% ==============================================================================

Building on the threefold division in Doc.~190 (P37) -- (1)~microscopic/$\xi$
level, (2)~macroscopic/$1/\varphi$ transit point, (3)~particle physics -- a
finer resolution is given here. Each role is provided with its sources and its
status.

\subsection{Role A -- Self-similarity attractor (asserted)}
\textbf{Source:} Doc.~172, \S\,``The golden ratio as a fractal attractor''.
\textbf{Statement:} The fractal field geometry with $D_f=3-\xi$ is said to
force a self-similarity recursion whose only stable fixed point is $\varphi$
(algebraically $\varphi^2=\varphi+1$).
\textbf{Status: asserted, not derived.} $\varphi^2=\varphi+1$ \emph{is} the
defining equation of $\varphi$; the actual step ``$D_f=3-\xi$ \emph{forces}
exactly this recursion'' is asserted, not shown. The variant in Doc.~188
(``$\varphi$ as the most stable incommensurable ratio'') is more strongly
motivated -- via the continued fraction $[1;1,1,\dots]$ $\varphi$ is the ``most
irrational'' number -- but likewise hinges on the geometry selecting precisely
this ratio.

\subsection{Role B -- CP phase $\arctan(\varphi^{-3})$ (genuine content)}
\textbf{Source:} Doc.~172, \S\,``Formal proof: $\arctan(\varphi^{-3})$ as a
geometric necessity''. \textbf{Statement:} The Fibonacci quotients
$F(n{+}1)/F(n)$ converge to $\varphi$ with errors of order
$\varphi^{-2n}/\sqrt5$; the geometric mean of first and second order is
$\sqrt{\varphi^{-2}\varphi^{-4}}=\varphi^{-3}=\sqrt5-2$. With the torsion
phase it follows that
$\delta = 270^\circ + \arctan(\varphi^{-3}) \approx 283.28^\circ$
(NOvA/T2K 2024: $283.1^\circ\pm14^\circ$).
\begin{correct}
\textbf{Status: the best-justified $\varphi$ usage in the corpus.} Here
$\varphi^{-3}$ is not a choice but the natural intermediate scale of Fibonacci
convergence; the exponent $3$ is forward-motivated via the Z3 order
($3$~generations/sectors). This is qualitatively different from a fitted
$\varphi$ power.
\end{correct}

\subsection{Role C -- Fine-structure constant $\alpha=\xi\,\varphi^3\,F(7)$}
\textbf{Source:} Doc.~172, \S\,``$\alpha$ from pure geometry''.
\textbf{Statement:}
\begin{equation}
  \alpha \approx \xi\cdot\varphi^3\cdot F(7) = \xi\cdot\varphi^3\cdot 13
  \;\Rightarrow\; \alpha^{-1} = 136.19 \quad(-0.62\%).
\end{equation}
\textbf{Status: numerically verified as correct,} but to be understood as one
of several $\alpha$ routes. The rigorous $\alpha$ derivation runs via
$\alpha=\xi E_0^2$ with $E_0=\sqrt{m_e m_\mu}$ and the power relation
$\alpha\propto\xi^{11/2}$ (Doc.~011); the $\varphi^3\cdot13$ form is the
geometrically ``clean'' variant, but it deviates from the measured value by
$0.62\%$.

\subsection{Role D -- Mass/generation ladder (heuristic, erroneous)}
\textbf{Sources:} Doc.~172 ($m_\mu/m_e\approx\varphi^{10}$),
Doc.~188 (lepton ladder $\varphi^5\pi^2$, $\varphi^5$), Doc.~009
($\varphi^{12}/(1-\xi)^2$), Doc.~022/035/133 ($\varphi^{\text{gen}}$ scaling).
\textbf{Status: heuristic; in Doc.~172/188 numerically wrong} (see
Section~\ref{sec:fehler}). The robust mass derivation runs \emph{not} via
$\varphi$ but via $\xi$ and quantum numbers (Doc.~006,
Section~\ref{sec:masse}). This role falls under P35.

\subsection{Role E -- Fractal scaling factor $\alpha=\varphi^{-1}$
(NOT an error)}
\textbf{Source:} Doc.~188, \S\,``Fractal coupling constant''.
\textbf{Statement:} $\alpha=\varphi^{-1}=0.6180$ as a scaling factor between
fractal levels, $(x_{k+1},\dots)=\varphi^{-1}(x_k,\dots)$.
\begin{important}
\textbf{Status: its own, clearly defined quantity -- not the fine-structure
constant.} In the same document the fine structure stands cleanly separated as
$\alpha_{\text{EM}}=\xi E_0^2=1/137.04$. Thus there is \emph{no} computational
error, but a deliberate double assignment of the letter $\alpha$. A pure
readability matter; Doc.~190 P37(a) already lists this usage.
\end{important}

\subsection{Role F -- Transit point $1/\varphi$ of the $\xi$ recursion}
\textbf{Source:} Doc.~275. \textbf{Statement:} The damping sequence
$r(k{+}1)=r(k)(1-\xi)$, $r(0)=D_f/3$, passes $1/\varphi$ at
$k^* = \log\varphi/\xi \approx 3609$ (exactly $k^*_{\text{exact}}=3608.51$).
\begin{caution}
\textbf{Status (revision 11 June 2026): transit point, not a fixed point.}
The only fixed point of the recursion is $0$. Since $k^*(c)$ exists for
\emph{any} target $c$, the observation is unanchored as long as $1/\varphi$ is
not independently motivated -- the $\varphi^N$ triviality warning P35 applies.
\end{caution}

\subsection{Role G -- Pentagonal symmetry breaking}
\textbf{Sources:} Doc.~018 (anomalous $g{-}2$), Doc.~149 (torsion),
Doc.~019, 153, 154. \textbf{Statement:} $\varphi$ as the carrier of pentagonal
(fivefold) symmetry breaking. \textbf{Status: structural motif;} it shares the
justification status of Role~A (the selection of precisely the pentagonal
symmetry is the open point).

\subsection{Role H -- HLV substrate (not FFGFT)}
\textbf{Source:} Doc.~271. \textbf{Statement:} In Hidden-Layer Vacuum (HLV,
Krueger) $\varphi$ is the substrate geometry. \textbf{Status: belongs to HLV,
not FFGFT.} In the bridge comparison $\varphi$ must not be tacitly transferred
from the HLV side to the FFGFT side.

\subsection{Role I -- Trivial representation $X=\varphi^N$}
\textbf{Sources:} Doc.~190 (P35), 271, 272, 274. \textbf{Statement and
status:} any magnitude can be written as $\varphi^N$ \emph{or} $\xi^N$; an
exponent alone is not a prediction. Holds equally for both bases.

% ==============================================================================
\section{Honest register: correct / open / wrong}\label{sec:fehler}
% ==============================================================================

\subsection{Correctly recomputed}
\begin{itemize}
  \item $\alpha = \xi\,\varphi^3\cdot13 \Rightarrow \alpha^{-1}=136.19$
        ($-0.62\%$), Doc.~172.
  \item $\varphi^{-3}=\sqrt5-2=0.23607$, $\arctan(\varphi^{-3})=13.28^\circ$,
        Doc.~172.
  \item $1/\varphi=0.6180$ as a transit point at $k^*\approx3609$,
        $r(3609)=0.617994\ldots$, Doc.~275 (status: unanchored).
  \item $245 = 5\cdot7^2 = \varphi^{12}/(1-\xi)^2 \approx 244.98$, Doc.~009 --
        with a \emph{declared} empirical anchor $1.314$ (P10/P37(d)).
\end{itemize}

\subsection{Numerically wrong -- in the legacy documents, corrected here}
\begin{critical}
\textbf{Doc.~188, lepton ladder (line~$\approx$318):}
\begin{align}
  \text{claimed: } & \frac{m_\mu}{m_e}\approx\varphi^5\cdot\pi^2\approx 206.7
    && \text{actually } \varphi^5\pi^2 = 109.46 \;(\text{measured }206.768),\\
  \text{claimed: } & \frac{m_\tau}{m_\mu}\approx\varphi^5\approx 16.94
    && \text{actually } \varphi^5 = 11.09 \;(\text{measured }16.82).
\end{align}
\textbf{Doc.~172 (line~$\approx$356):} $m_\mu/m_e\approx\varphi^{10}=122.99$ --
$40\%$ off and contradicting the (likewise wrong) Doc.~188 form. Even
arithmetically corrected one would need awkward exponents
($\varphi^{11.08}$ resp.\ $\varphi^{5.87}$): there is no clean
$\varphi$ power that hits the lepton ratios.
\textbf{Consequence:} The $\varphi$ lepton ladder is to be flagged as
numerology (P35 null result), not as a geometric consequence.
\end{critical}

\subsection{The robust lepton-mass derivation -- via $\xi$, not
$\varphi$}\label{sec:masse}
Doc.~006 sets $m_i = r_i\,\xi^{p_i}\,v$ with exponents fixed from quantum
numbers:
\begin{equation}
  p_e = \tfrac32,\; r_e\approx 1.349; \qquad
  p_\mu = 1,\; r_\mu = \tfrac{16}{5} = 3.2;
\end{equation}
\begin{equation}
  \frac{m_\mu}{m_e} = \frac{r_\mu}{r_e}\,\xi^{p_\mu-p_e}
  = 2.372\cdot\xi^{-1/2} = 2.372\cdot86.6 = 205.4 \quad(-0.65\%).
\end{equation}
The \emph{exponent difference} $-\tfrac12$ is the structural, forward-directed
content; the $O(1)$ coefficients ($r_e$ is back-computed) carry the residual
fit (cf.\ P40: only ratios exact). $\varphi$ is not needed here and adds
nothing.

% ==============================================================================
\section{What does NOT follow from $\varphi$ (or $\xi$)}
% ==============================================================================

\begin{itemize}
  \item \textbf{$\pi$ is not derived.} $\pi$ is the presupposed circle/torus
        constant. The nicest interpretation (Doc.~159): on the circle
        $\sum 1/n^2=\pi^2/6$ converges, the periodic geometry forces
        convergence -- but $\pi$ is interpreted, not derived.
  \item \textbf{$e$ (Euler's number) is not derived}, it appears only as
        $\exp(\cdot)$.
  \item \textbf{Emergently derived}, by contrast, are: $13=F(7)$ (Fibonacci)
        and $27=3^3$ (Z3 structure, three generations).
\end{itemize}

% ==============================================================================
\section{Table of occurrences (substantial cases)}
% ==============================================================================

The table lists the \emph{content-relevant} occurrences. Pure definitional or
pass-through mentions follow the same canonical reading.

\begin{center}
\begin{adjustbox}{max width=\textwidth}
\begin{tabular}{l l l l}
\hline
\textbf{Doc.} & \textbf{Symbol} & \textbf{Role (see above)} & \textbf{Status} \\
\hline
172 & $\varphi$ & A attractor, B CP phase, C $\alpha$, D $\varphi^{10}$ & B correct; A asserted; D wrong \\
188 & $\phi$ & D lepton ladder, E scaling factor & D wrong; E ok (own quantity) \\
275 & $\varphi$ & F transit point $1/\varphi$ & correct, but unanchored (P35) \\
190 & $\varphi$ & I triviality P35, role map P37 & normative \\
009 & $\phi$ & D $\varphi^{12}/1.314$ & empirical anchor declared \\
271 & $\varphi$ & H HLV substrate, I triviality & HLV/FFGFT delineation \\
272 & $\varphi$ & I $X=\varphi^N$ not a prediction & correct \\
274 & $\varphi$ & I triviality holds for $\varphi$ and $\xi$ & correct \\
018 & $\varphi$ & G pentagonal symmetry breaking ($g{-}2$) & structural motif \\
149 & $\varphi$ & G pentagonal symmetry breaking (torsion) & structural motif \\
019, 153, 154 & $\varphi$ & G pentagonal symmetry breaking & structural motif \\
022, 035, 133 & $\phi$ & D $\varphi^{\text{gen}}$ scaling & heuristic (P35) \\
036 & $\phi$ & D cosmic scales $r_n=r_0\varphi^n$ (BAO) & heuristic (P35) \\
083, 147, 187 & $\phi$ & D qubit/photonics frequencies & heuristic \\
005, 028, 035, 150 & $\phi$/$\varphi$ & definition $(1+\sqrt5)/2$ & basis \\
241, 261 & $\varphi$ & F ``settled'' recursion & cf.\ Doc.~275 \\
\hline
202 & $\phi$ & \emph{scalar field} $\phi_{n,l,j}$, index $n_\phi$ & NOT golden ratio \\
051 & $\phi$ & \emph{winding index} $n_\phi$ & NOT golden ratio \\
067, 097, 129 & $\phi$ & \emph{scalar field} (Lagrangian/QFT) & NOT golden ratio \\
\hline
\end{tabular}
\end{adjustbox}
\end{center}

% ==============================================================================
\section{Recommendation}
% ==============================================================================

\begin{enumerate}
  \item \textbf{Notation going forward:} golden ratio exclusively
        \verb|\varphi| ($\varphi$). \verb|\phi| remains field/angle,
        \verb|\Phi| potential, \verb|\psi| wave function.
  \item \textbf{Legacy documents:} do not rewrite; this reference governs.
  \item \textbf{$\varphi$ mass ladders (Doc.~172/188):} read as P35 numerology,
        not as derivation. Only the $\xi$ route is robust (Doc.~006).
  \item \textbf{Genuine $\varphi$ content} remains: $\arctan(\varphi^{-3})$ in
        the CP phase (Doc.~172) and the conceptual separation of the three
        $\varphi$ roles (Doc.~190 P37).
\end{enumerate}
